\documentclass[11pt]{article}
\usepackage[a4paper,margin=2.5cm]{geometry}
\usepackage{hyperref}
\usepackage{xcolor}
\usepackage{listings}
\usepackage{parskip}
\usepackage{titlesec}
\usepackage{enumitem}

\hypersetup{colorlinks=true, linkcolor=blue, urlcolor=blue}
\titleformat{\section}{\Large\bfseries}{\thesection}{1em}{}
\titleformat{\subsection}{\large\bfseries}{\thesubsection}{1em}{}

\definecolor{codebg}{RGB}{245,245,245}
\lstset{
  backgroundcolor=\color{codebg},
  basicstyle=\ttfamily\small,
  breaklines=true,
  frame=single,
  columns=fullflexible,
  keepspaces=true
}

\title{\textbf{MIMaaS + Hailo-10H Setup Manual}\\
\large Raspberry Pi 5, Hailo Dataflow Compiler, and the MIMaaS server}
\author{Maxime \\ AnkiLab}
\date{August 2026}

\begin{document}
\maketitle
\tableofcontents
\newpage

\section{Overview}

This manual describes how to set up the hardware and software required to
run the Hailo-10H backend of MIMaaS end-to-end, from a fresh Raspberry Pi~5
to a successful \texttt{flash\_and\_evaluate()} call. It assumes the reader
already has the base MIMaaS server installed (see the project
\texttt{README.md} for the generic installation of the Python package,
Redis, and the database); this document focuses \emph{specifically} on the
pieces needed for the Hailo-10H target: the Pi~5 itself, the Hailo Dataflow
Compiler (DFC), and the WSL2 bridge used to invoke it from a Windows
development machine.

\textbf{Architecture recap.} Model compilation (TFLite $\rightarrow$ HEF)
happens on the \emph{development/server machine} inside WSL2, using the
Hailo DFC toolchain. The resulting \texttt{.hef} binary is then transferred
over SSH to the Raspberry~Pi~5, which hosts the physical Hailo-10H module,
and executed there via \texttt{hailortcli}. The server machine never needs
direct access to the NPU; the Pi only ever needs \texttt{hailortcli} and the
runtime, not the full DFC.

\section{Part A -- Raspberry Pi 5 + Hailo-10H}

\subsection{Hardware}
\begin{itemize}
  \item Raspberry Pi 5 (4GB or 8GB)
  \item Hailo-10H M.2/HAT accessory, correctly seated and powered
  \item Network connectivity between the Pi and the machine running MIMaaS
    (in this deployment, a static link at \texttt{192.168.42.1})
\end{itemize}

\subsection{Software on the Pi}
\begin{enumerate}
  \item Flash a recent Raspberry Pi OS (64-bit) image and boot the Pi.
  \item Install the HailoRT runtime and \texttt{hailortcli} following
    Hailo's official Raspberry Pi installation guide for the Hailo-10H
    (PCIe/driver setup, \texttt{hailo-pci} kernel module, HailoRT service).
    Verify the module is detected:
\begin{lstlisting}
hailortcli fw-control identify
\end{lstlisting}
  This should report the Hailo-10H device and firmware version. If it
  reports "no device found", check the physical connection and that the
  \texttt{hailo\_pci} kernel module is loaded (\texttt{lsmod | grep hailo}).
  \item Enable passwordless SSH key authentication from the MIMaaS server
    machine to the Pi (the backend uses \texttt{BatchMode=yes}, i.e.\ it
    will \emph{not} prompt for a password):
\begin{lstlisting}
# On the server machine (or inside WSL2):
ssh-keygen -t ed25519 -f ~/.ssh/mimaas_hailo_pi5
ssh-copy-id -i ~/.ssh/mimaas_hailo_pi5.pub maxime@192.168.42.1

# Verify non-interactive login works:
ssh -o BatchMode=yes maxime@192.168.42.1 echo ok
\end{lstlisting}
  \item Create a working directory on the Pi for uploaded models, matching
    the path convention used by the backend (\texttt{/tmp/mimaas/<board>/}
    for execution, and any directory of your choosing for the source
    \texttt{.hef} files, e.g.\ \texttt{/home/maxime/models/}).
\end{enumerate}

\section{Part B -- Hailo Dataflow Compiler (DFC) in WSL2}

The DFC (compiler, optimizer, parser) currently runs only on x86\_64 Linux.
On a Windows development machine, install it inside WSL2.

\subsection{Install WSL2 and a pyenv-managed Python 3.10}
\begin{lstlisting}
wsl --install -d Ubuntu-22.04

# Inside the WSL2 Ubuntu shell:
curl https://pyenv.run | bash
# add pyenv init to ~/.bashrc as instructed by the installer, then:
pyenv install 3.10.14
pyenv virtualenv 3.10.14 hailo_dfc_530_py310
pyenv activate hailo_dfc_530_py310
\end{lstlisting}

\subsection{Install Hailo DFC 5.3.0}
Download the DFC \texttt{.whl} package for your DFC version from the Hailo
Developer Zone (requires a free Hailo account) and install it inside the
virtual environment created above:
\begin{lstlisting}
source ~/.pyenv/versions/3.10.14/envs/hailo_dfc_530_py310/bin/activate
pip install hailo_dataflow_compiler-5.3.0-*.whl
hailo --version   # sanity check
\end{lstlisting}

\textbf{Note on GPU}: without a supported NVIDIA GPU visible inside WSL2,
DFC automatically falls back to CPU-only optimization at level~0
(``the accuracy won't be optimized and compression won't be used''). This is
sufficient for compatibility/performance testing (this manual and the
accompanying report's results were produced this way) but is \textbf{not}
representative of production accuracy -- for production deployments, run
optimization on a machine with a compatible GPU and a real calibration
dataset (not \texttt{--use-random-calib-set}).

\subsection{Environment variables used by the MIMaaS backend}
The backend (\texttt{src/mimaas/hardware/hailo\_utils.py}) picks up the
following environment variables, with sensible defaults:
\begin{center}
\begin{tabular}{ll}
\textbf{Variable} & \textbf{Default} \\
\hline
\texttt{HAILO\_DFC\_VENV} & \texttt{hailo\_dfc\_530\_py310} \\
\texttt{HAILO\_HW\_ARCH} & \texttt{hailo10h} \\
\texttt{HAILO\_SSH\_PREP\_TIMEOUT\_S} & 20 \\
\texttt{HAILO\_SCP\_TRANSFER\_TIMEOUT\_S} & 120 \\
\texttt{HAILO\_SSH\_INFERENCE\_TIMEOUT\_S} & 900 \\
\end{tabular}
\end{center}
If your pyenv virtualenv has a different name, set
\texttt{HAILO\_DFC\_VENV} accordingly before starting the MIMaaS worker.

\section{Part C -- MIMaaS Server Configuration}

\subsection{Board configuration file}
Add or verify the board definition at
\texttt{config/boards/hailo\_pi5.yaml}:
\begin{lstlisting}
value:
- variant: hailo_rpi5
  board_type: hailo_rpi5
  board_snr: "HAILO10_PI5_01"
  ppk2_snr: "N/A"
  uart_port_type: none
  voltage: 5000
  power_measurement_threshold: 1000
  max_available_tensor_arena_size: 0
  flash_size: 32768
  ram_size: 2048
  host: "192.168.42.1"       # <- IP of your Raspberry Pi 5
  ssh_user: "maxime"          # <- SSH user on the Pi
  processor: "Cortex-A76"
  npu: "Hailo-10"
  supported_formats:
    - ".hef"
    - ".tflite"
  max_model_size_mb: 50
\end{lstlisting}
Adjust \texttt{host} and \texttt{ssh\_user} to match your Pi's network
configuration and SSH account.

\subsection{Install the MIMaaS package}
On the machine that will run the API/worker (can be Windows; the DFC calls
are automatically routed into WSL2 by the backend):
\begin{lstlisting}
git clone <repository-url> mimaas-server
cd mimaas-server
python3 -m venv .venv
# Windows:
.venv\Scripts\activate
# Linux/macOS:
source .venv/bin/activate

pip install -e .
\end{lstlisting}

\section{Part D -- Running a Full Test}

\subsection{Ad-hoc test via the pipeline API}
The fastest way to validate an end-to-end setup is to call
\texttt{flash\_and\_evaluate()} directly:
\begin{lstlisting}
from mimaas.core.pipeline import flash_and_evaluate
from mimaas.models.board import Board

board = Board.from_config("hailo_pi5")

# Case 1: a pre-compiled .hef -- skips DFC entirely
result = flash_and_evaluate("/path/to/model.hef", board)

# Case 2: a .tflite -- triggers parser -> optimize -> compiler in WSL2
result = flash_and_evaluate("/path/to/model.tflite", board)

print(result)
# {'status': 'flash_and_eva_successful', ...}
\end{lstlisting}

A successful result includes the parsed \texttt{hailortcli} summary: FPS and
Hailo-10H chip temperature (mean/min/max), in addition to the standard
MIMaaS result fields (ROM usage in MB, host RAM usage).

\subsection{Via the REST API / worker (production path)}
Once the API server and a Celery worker for the \texttt{hailo\_pi5} queue are
running (see \texttt{docs/START\_SYSTEM.md} for the generic API/worker/Redis
bring-up steps), submit a request exactly as for MCU boards:
\begin{lstlisting}
curl -X POST http://127.0.0.1:5000/api/requests/ \
  -H "Authorization: YOUR_API_TOKEN" \
  -F "network=@/path/to/model.tflite" \
  -F "board=hailo_pi5"
\end{lstlisting}

\section{Troubleshooting}

\subsection{``FileNotFoundError: efficientnetb3\_opt.har'' (or similar) after
an \texttt{optimize} failure}
If \texttt{hailo optimize} fails partway through (e.g.\ a quantization
error), it does \emph{not} write the output \texttt{\_opt.har} file. Do not
chain the \texttt{compiler} command unconditionally after
\texttt{optimize} in a shell script -- check the optimize step's exit code
first, or run steps one at a time as shown in this manual.

\subsection{``Shift delta in .../avgpool.../avgpool\_op is larger than 2''}
This is a quantization limitation of the model's global-average-pool layer(s),
not a bug in your setup. Try forcing that specific layer to 16-bit precision
via a model script before re-running \texttt{optimize}:
\begin{lstlisting}
cat > model.alls <<'ALLS'
quantization_param(avgpool1, precision_mode=a16_w16)
ALLS

hailo optimize model.har --hw-arch hailo10h --use-random-calib-set \
  --model-script model.alls --output-har-path model_opt.har
\end{lstlisting}
If the error reappears on a \emph{different} average-pool layer name (e.g.\
\texttt{avgpool2}), the model likely has multiple squeeze-and-excitation
blocks and this single-layer workaround will not scale; see Section~5.2.1 of
the accompanying report for a discussion of when this is/is not fixable.

\subsection{Parser interactively asks ``Would you like to parse again with
the recommendation? (y/n)'' and then crashes with \texttt{EOFError}}
This happens when the parser fails and offers an automatic recommendation,
but is invoked without a TTY (e.g.\ from a Python subprocess), so it cannot
read the interactive answer. Re-run the \texttt{hailo parser} command with
the \texttt{-y} flag to auto-accept parser recommendations. Note that for
some architectures (e.g.\ ConvNeXt, whose LayerNorm decomposition uses
operations the parser does not support at all) accepting the recommendation
still will not produce a working graph -- see the failure taxonomy in the
report.

\subsection{WSL2 runs out of memory during compilation}
Large models (100M+ parameters, e.g.\ VGG16) can exhaust a WSL2 VM's default
memory allocation. Monitor memory during a compile with
\texttt{watch -n 2 free -h} in a second WSL2 terminal, and avoid running more
than one DFC compilation at a time. If needed, increase the WSL2 memory cap
in \texttt{.wslconfig} (Windows user profile directory) and restart WSL2
(\texttt{wsl --shutdown}).

\subsection{SSH commands hang or time out}
Confirm the Pi is reachable (\texttt{ping 192.168.42.1}) and that
\texttt{ssh -o BatchMode=yes maxime@192.168.42.1 echo ok} succeeds
non-interactively \emph{before} involving MIMaaS -- \texttt{BatchMode=yes}
means any authentication prompt will fail the connection immediately rather
than hang.

\end{document}
